\section{Institutional Policy Recommendations}
\label{sec:concl_recommendations}

The empirical findings translate into four actionable institutional directives for higher education administrators:

\begin{enumerate}
    \item \textbf{Transition from Correlational Scoring to Counterfactual Simulation:} Deprecate predictive "early warning" risk scores that only report dropout probability. Deploy causal graphical engines capable of evaluating counterfactual interventions ($P(Y \mid do(T=t), \mathbf{x})$) before public funds are committed.
    \item \textbf{Overhaul Merit-Only Allocation (Targeting Vulnerability):} Shift away from allocating aid solely on entry grades ($Z_2$), which subsidizes affluent students who would graduate regardless. Prioritize non-traditional older students ($>20$ years) and financially indebted students, who exhibit treatment multipliers of $1.72\times$ and $2.03\times$.
    \item \textbf{Couple Financial Aid with Early Academic Tutoring:} Because $21.53\%$ of the scholarship's causal effect is mediated through first-semester course completions ($T \to M_1 \to Y$), financial aid must be paired with proactive coaching during Semester 1 to protect academic momentum.
    \item \textbf{Strict Intervention Window (End of Semester 1):} By the end of Semester 1, course approval records provide definitive identification of students in the dropout spiral. Intervening after Semester 2 is largely ineffective because student departure has already crystallized.
\end{enumerate}
